\section{Introduction}
\label{sec:intro}
A visual policy can propose several plausible action chunks, yet small differences between them can decide whether an object reaches its target. Sampling additional proposals helps only when a selector can identify the useful ones. This makes the prediction target a practical design choice: what consequence should a world model forecast so that a reader can rank actions accurately and economically?

Existing approaches offer different answers. Visual world models predict future spatial features or observations~\cite{zhou2025dinowm,assran2025vjepa2}; these targets support image-goal evaluation but retain substantial spatial structure. Compact visual tokenizers reduce that representation~\cite{kim2026compact}. Direct value guidance instead scores candidate actions~\cite{nakamoto2024vgps}, while predictive policy steering evaluates their generated consequences~\cite{qi2025gpc,cui2026dreamsteer}. These methods establish that prediction, compression, and proposal selection can each benefit control. They leave a focused empirical question: how useful is a learned, context-conditioned \emph{segment target} between a full visual forecast and a direct score, when proposals and evaluation states are held fixed?

We study this question with Conditional Trajectory Abstraction (CTA). During training, an encoder observes context $C$ and sampled future observations $\tau$, producing a compact code $S$. A predictor learns $p(S\mid C,\bA)$ for a proposed action chunk $\bA$. At deployment, a reader evaluates the prediction together with $C$ and a goal image $g$. Its only candidate-specific input is the predicted code: it cannot bypass the representation by reading $\bA$. Actual future observations are training targets, never deployment inputs. Context conditioning permits the source and reader to share information, and segment supervision includes intermediate observations; the benefits of those individual choices still require isolated ablations.

This design is related to conditional trajectory codes~\cite{jiang2023tap}, context-aware tokenization~\cite{micheli2024deltairis}, and decision-aware model compression~\cite{grimm2020value,arumugam2022deciding}. Our methodological contribution is their concrete use as a future-supervised target for visual action ranking: the code is learned from observed consequences, forecast from a supplied action chunk, and evaluated through an action-blind reader. The contribution concerns this target-learning interface and its decision quality, building on those established ingredients.

We ask three questions. \textbf{RQ1:} Does this target improve ranking on identical candidate banks relative to endpoint prediction and direct scoring? \textbf{RQ2:} Does ranking quality translate into closed-loop success and useful computation savings? \textbf{RQ3:} Which failures arise in source encoding, prediction, or the evaluation protocol? The answers are deliberately separated. CTA improves development-set geometry ranking, including against a larger direct scorer. Its historical closed-loop advantage over learned alternatives is much less certain. Component timing shows a compact scoring path, but policy sampling dominates and total-planner savings are not established. Longer-horizon prediction errors and sensitivity to proposal batch shape explain concrete limits of the current evidence.

The paper contributes: (1) a conditional segment-target architecture with an explicit deployment information boundary; (2) a two-stage training recipe combining future-feature reconstruction, candidate ranking, and forecast-to-reader consistency; and (3) a development evaluation separating privileged future reading, predicted-code ranking, historical closed-loop behavior, and numerical reproducibility. This third contribution includes negative findings: a lower prediction loss can accompany code collapse, and a ranking gain need not produce a reliable episode-level gain. The goal is to establish what this representation currently buys and what evidence is still needed for a general quality--compute claim.

\begin{figure*}[t]
\centering
\resizebox{\linewidth}{!}{%
\begin{tikzpicture}[x=1cm,y=1cm,>=Latex,
 box/.style={draw=black!55,rounded corners=2pt,fill=white,minimum height=.70cm,align=center,font=\sffamily\small},
 flow/.style={->,line width=.8pt}, aux/.style={->,line width=.7pt,draw=black!55},
 note/.style={font=\sffamily\footnotesize,align=center}]
\path[use as bounding box] (0,0) rectangle (17.3,6.15);
\node[note,anchor=west,font=\sffamily\bfseries\small] at (.1,5.91) {Stored PushT inputs};
\foreach \xx/\file/\lab in {2.8/prev/Previous,4.2/cur/Current,7.1/c3s2/Step 2,8.5/c3s4/Step 4,9.9/c3s6/Step 6,11.3/c3s8/Step 8,15.0/goal0/Goal} {
 \node[inner sep=0] at (\xx,5.13) {\includegraphics[width=1.02cm]{fig/pusht/\file.png}};
 \node[note] at (\xx,4.47) {\lab};
}
\draw[black!20] (.1,4.18)--(17.2,4.18);
\node[note,anchor=west,font=\sffamily\bfseries\small] at (.1,3.90) {TRAIN: actual futures available};
\node[box,minimum width=2.45cm] (ct) at (1.6,3.12) {Observed context $C$};
\node[box,minimum width=2.45cm,fill=orange!9] (ft) at (1.6,2.15) {Sampled future $\tau$};
\node[box,minimum width=2.1cm] (enc) at (5.1,2.63) {Source encoder\\$q_\phi(C,\tau)$};
\node[box,minimum width=1.6cm,fill=blue!9] (code) at (8.0,2.63) {FSQ target $S$\\16 tokens};
\node[box,minimum width=2.1cm] (read) at (11.0,3.18) {Reader $D_\psi$\\$C,S,g$};
\node[box,minimum width=2.1cm] (dec) at (11.0,2.14) {Feature decoder\\$C,S$};
\node[note,anchor=west] (rank) at (13.0,3.18) {Sibling ranking supervision};
\node[note,anchor=west] (rec) at (13.0,2.14) {Future image-feature reconstruction};
\draw[flow] (ct.east)--(enc.north west);
\draw[flow] (ft.east)--(enc.south west);
\draw[flow] (enc)--(code);
\draw[flow] (code.east)--(read.west);
\draw[flow] (code.east)--(dec.west);
\draw[aux] (read)--(rank);
\draw[aux] (dec)--(rec);
\draw[black!20] (.1,1.55)--(17.2,1.55);
\node[note,anchor=west,font=\sffamily\bfseries\small] at (.1,1.32) {DEPLOY: actual futures unavailable};
\node[box,minimum width=2.7cm] (inp) at (1.7,.51) {Context $C$ + proposal $\bA$};
\node[box,minimum width=2.1cm] (wm) at (5.1,.51) {Predictor $h_\theta$};
\node[box,minimum width=2.1cm,fill=blue!9] (mean) at (8.0,.51) {Forecast $\hat S$\\48 continuous scalars};
\node[box,minimum width=2.1cm] (dr) at (11.0,.51) {Frozen reader\\$C,\hat S,g$};
\node[box,minimum width=2.2cm] (sel) at (15.0,.51) {Rank, execute, replan};
\draw[flow] (inp)--(wm);
\draw[flow] (wm)--(mean);
\draw[flow] (mean)--(dr);
\draw[flow] (dr)--node[note,above] {score} (sel);
\node[note,anchor=south] at (11.0,.90) {shared context + goal features};
\end{tikzpicture}}

\caption{\textbf{CTA learns what the predictor must forecast.} Training uses observed context and sampled future observations to construct a quantized segment target. Ranking and feature reconstruction shape that target; the frozen reader then guides forecast training. Deployment uses context, policy proposals, and goal features. Only the predictor receives candidate actions; only training receives actual futures. The reader receives a continuous predictive mean. Images illustrate stored PushT observations, not a claimed successful deployment.}
\label{fig:overview}
\end{figure*}
